\documentclass[a4paper,10pt]{article}
\usepackage[italian]{babel}
% Shared A4 layout. Each language file selects its own Babel language.
\usepackage[margin=16mm]{geometry}
\usepackage[T1]{fontenc}
\usepackage[utf8]{inputenc}
\usepackage{lmodern}
\usepackage{xcolor}
\usepackage{enumitem}
\usepackage{titlesec}
\usepackage{needspace}
\usepackage{fancyhdr}
\usepackage{url}

\definecolor{accent}{HTML}{29634D}
\definecolor{muted}{HTML}{606964}
\pdfinfo{
  /Title (Riccardo Giorgi - Curriculum Vitae)
  /Author (Riccardo Giorgi)
}
% url prints readable addresses without creating PDF hyperlink annotations.
\urlstyle{same}
\setlength{\parindent}{0pt}
\setlength{\parskip}{0.25em}
\setlength{\emergencystretch}{2em}
\setlist[itemize]{leftmargin=1.2em, nosep, topsep=0.3em}
\titleformat{\section}{\large\bfseries\color{accent}}{}{0pt}{}[\titlerule]
\titlespacing*{\section}{0pt}{1em}{0.5em}
\pagestyle{fancy}
\fancyhf{}
\renewcommand{\headrulewidth}{0pt}
\fancyfoot[L]{\small\color{muted}Riccardo Giorgi}
\fancyfoot[R]{\small\color{muted}\thepage}

% Title, dates, organisation/location, description.
\newcommand{\cventry}[4]{%
  \Needspace{4\baselineskip}%
  \textbf{#1}\par
  {\small\color{muted}#2\enspace\textbullet\enspace#3}\par
  #4\par\vspace{0.2em}
}

\newenvironment{cvreferences}{%
  \begingroup\small
  \begin{description}[leftmargin=0pt, style=nextline, itemsep=0.1em, parsep=0pt]
}{%
  \end{description}
  \endgroup
}


\begin{document}

{\Huge\bfseries Riccardo Giorgi}\par
{\large Studente di Ingegneria Meccanica}\par
{\color{muted}Politecnico di Torino}\par
riccardogiorgi29@gmail.com
\enspace\textbullet\enspace +39 371 387 0042
\enspace\textbullet\enspace 10141 Torino, Italia\par
\url{riccardogiorgi.tech/it/}\par
{\small Patenti di guida: AM, B\enspace\textbullet\enspace Partita IVA: 13339980016}

\section*{Profilo}
Studio Ingegneria Meccanica al Politecnico di Torino, nel corso erogato in
inglese, e cerco un tirocinio in ambito meccanico o strutturale, preferibilmente
a Torino. Uso la programmazione e la modellazione 3D a supporto della
progettazione, della simulazione e della risoluzione di problemi tecnici.
Imparo rapidamente e lavoro bene nei gruppi che danno valore alla logica,
al metodo e a una comunicazione chiara.

\section*{Esperienze}
\cventry{Membro dei gruppi Strutture e Firmware}{Ottobre 2025 - oggi}
  {Team Icarus - Advanced Propulsion Aircraft, Politecnico di Torino}
  {Progettazione e analisi strutturale con SolidWorks, HyperMesh e HyperView:
  preparazione dei modelli, generazione della mesh, carichi e analisi dei
  risultati. Da ottobre 2026 contribuisco anche al firmware di controllo del
  volo e telemetria in C su microcontrollori STM32 con FreeRTOS e MAVLink,
  inclusa la comunicazione UART e I\textsuperscript{2}C.}

\cventry{Sviluppatore software part-time}{Dicembre 2023 - ottobre 2025}
  {TBJBU2 Ltd, A distanza}
  {Ho contribuito a diversi titoli commerciali attraverso implementazione,
  integrazione e affidabilità in produzione, mettendo a frutto l'esperienza
  dei progetti personali in C++ e Unreal Engine. I dettagli sono coperti
  da accordi di riservatezza.}

\cventry{Tutor di programmazione}{2023}{Kodland, A distanza}
  {Ho insegnato le basi della programmazione a classi di circa dieci bambini,
  spiegando concetti tecnici ad allievi e genitori e sviluppando capacità di
  comunicazione e gestione del tempo.}

\cventry{Cameriere}{Luglio 2022}{Monica Li, Infernetto, Roma}
  {Ho lavorato in un ristorante con ritmi sostenuti, sviluppando capacità di
  adattamento, coordinazione e lavoro di squadra.}

\section*{Formazione universitaria}
\cventry{Laurea in Ingegneria Meccanica - in corso}
  {Settembre 2024 - oggi}{Politecnico di Torino}
  {Corso erogato in inglese, con l'ingegneria meccanica come principale
  obiettivo accademico e professionale.}

\cventry{Percorso di eccellenza INTRAPRENDENTI}{Settembre 2024 - oggi}
  {Politecnico di Torino}
  {Percorso interdisciplinare per studenti meritevoli, frequentato insieme
  alla laurea, con laboratori, gruppi multidisciplinari e sviluppo delle
  competenze trasversali.}

\newpage

\section*{Qualifiche e certificati}
\cventry{Badge Grandi Sfide Globali}{22 luglio 2026}{Politecnico di Torino}
  {Risoluzione multidisciplinare dei problemi, pensiero critico,
  progettazione responsabile e collaborazione tra discipline STEM,
  scienze umane e sociali.}

\cventry{Corso di Car Design - certificazione Autodesk Maya}
  {Completato il 17 aprile 2025}{Istituto Volta, A distanza}
  {Concetti di progettazione automobilistica e tecniche avanzate di
  modellazione 3D; credenziale verificabile su Credly.}

\cventry{International Baccalaureate}{Luglio 2023}
  {Southlands British International School, Roma}{Voto finale: 38/45.}

\cventry{Edexcel IGCSE}{Luglio 2021}
  {Southlands British International School, Roma}{Voto finale: 49/63.}

\section*{Progetto personale in corso}
\textbf{Videogioco di sopravvivenza e automazione}\par
Progettazione e sviluppo individuale di un videogioco in prima persona in
Unreal Engine e C++, per collegare raccolta, creazione di oggetti, ricerca e
prime forme di automazione in un prototipo giocabile, con salvataggi affidabili
e test con giocatori esterni come obiettivi.

\section*{Competenze e altre informazioni}
\begin{itemize}
  \item \textbf{Ingegneria:} SolidWorks; analisi FEM con HyperMesh e HyperView; Python; MATLAB.
  \item \textbf{Programmazione:} C++; Unreal Engine; C/STM32; FreeRTOS; MAVLink; Git; CMake.
  \item \textbf{Modellazione 3D:} Autodesk Maya (certificato); Blender (oltre 250 ore).
  \item \textbf{Lingue:} italiano madrelingua; inglese C2; francese B2.
  \item \textbf{Mobilità:} disponibile per tirocini e trasferte in Italia e nell'UE; preferenza per Torino.
  \item \textbf{Interessi:} scacchi rapid (10+0), massimi di 1709 su Chess.com e 1913 su Lichess; teatro (2018-2022).
\end{itemize}

\section*{Profili, credenziali e portfolio}
\begin{cvreferences}
  \item[LinkedIn] \url{linkedin.com/in/riccardo-giorgi-67b1b2240/}
  \item[GitHub] \url{github.com/Theuntextured}
  \item[Team Icarus - strutture e firmware]
    \url{icarus.polito.it/apa-2025/} \quad \url{icarus.polito.it/it-2025/}
  \item[Certificazione Maya] \url{credly.com/badges/286b9940-d45e-47c7-855e-a76377268bc5/}
  \item[Badge Grandi Sfide Globali] \url{bestr.it/award/show/NZdQ4dvER7mbhSBRxjHZKA}
  \item[FAB - pubblicazioni per Unreal Engine] \url{fab.com/sellers/The%20untextured%20Dev}
  \item[Unintended Upgrades - game jam] \url{the-untextured.itch.io/unintended-upgrades}
  \item[Scacchi online] \url{chess.com/member/the_untextured} \quad \url{lichess.org/@/The-untextured}
\end{cvreferences}

\end{document}
